% INSERCIÓN FOCAL AUTORIZADA: no introduce un capítulo ni el concepto de radión.
% Incluir dentro de la sección constitutiva, después de la inversión cúbica
% y de la declaración de las cartas dimensionales. Etiquetas fuera de grupos
% de prefijado automático. Conserva íntegramente el bloque geométrico anterior.
\subsection{Reconstruction de la forme d’action à partir de la réponse constitutive}
\label{iii:subsec:puente-forma-LC}

La réponse constitutive et la réalisation inductive--capacitive reçoivent
le même couple angulaire construit à partir d’APP--TRIT--TPK. La première évalue
des fonctions rationnelles des deux canaux ; la seconde utilise le rapport de
leurs exposants pour déterminer la forme de l’ellipse d’action.
L’inversion cubique permet de composer les deux lectures sans identifier leurs
quotients numériques.

Soient $x=A_{\rm rad}$ et $y=C^*_{\rm rad}$, avec $x>|y|$, les coordonnées
relevées déjà construites, et conservons leurs étiquettes de feuillet. Nous écrivons
\begin{equation}
 q_\pm=e^{-(x\pm y)},\qquad s_\pm=q_\pm^{30},\qquad
 r_\pm=f(s_\pm),\qquad
 f(s)=\frac{1+s+s^2}{1+s+s^2+s^3}.
 \label{iii:eq:forma-canales}
\end{equation}
Le domaine contractant donne $s_\pm\in(0,1)$ et $r_\pm\in(3/4,1)$. Dans la chambre $x>y>0$, on a $s_+<s_-$ et $r_+>r_-$ ; l'échange des feuillets transporte cette chambre vers l'orientation opposée.

L’image des deux coordonnées constitutives normalisées est le domaine
\begin{equation}
 \mathcal U=\left\{(z,v)\in(0,\infty)^2:
   (zv)^{-1/2}\in(3/4,1),\quad
   (z/v)^{1/2}\in(3/4,1)\right\}.
 \label{iii:eq:forma-dominio}
\end{equation}
Ici, $(z,v)=(\widehat Z,\widehat c)$. Les restrictions expriment l'appartenance des canaux reconstruits au domaine du théorème d'inversion ; elles ne constituent pas un choix ultérieur de leurs valeurs.

\Needspace{23\baselineskip}
\begin{proposition}[Composition constitutive de la forme et de l’impédance LC]
\label{iii:prop:forma-LC}
Pour $(\widehat Z,\widehat c)\in\mathcal U$, soient $s_\pm\in(0,1)$ les
racines uniques de
\begin{equation}
 r_\pm s_\pm^3+(r_\pm-1)(1+s_\pm+s_\pm^2)=0,
 \qquad
 r_+=\frac1{\sqrt{\widehat Z\widehat c}},\quad
 r_-=\sqrt{\frac{\widehat Z}{\widehat c}}.
 \label{iii:eq:forma-inversion}
\end{equation}
Le rapport angulaire, l’anisotropie de l’ellipse d’action et l’impédance
LC relative sont reconstruits exactement par
\begin{align}
 \frac{y}{x}
 &=\frac{\log s_+-\log s_-}{\log s_++\log s_-},
 \label{iii:eq:forma-razon}\\
 \eta_{\rm el}
 &=\frac12\log\frac{\log s_+}{\log s_-},
 \label{iii:eq:forma-eta}\\
 \frac{Z_{\rm LC}}{Z_*}
 &=e^{-\eta_{\rm el}}
  =\sqrt{\frac{\log s_-}{\log s_+}}.
 \label{iii:eq:forma-LC}
\end{align}
Dans la dernière identité, $Z_{\rm LC}=Z_{\eta_{\rm el}}$ est l’impédance
de la coordinatisation LC de \eqref{hmtII:eq:ii-elipse-lc}.
\end{proposition}
\begin{proof}
L'inversion cubique de \eqref{iii:eq:cubic} détermine $s_\pm$ de manière univoque, car $f$ est une bijection décroissante de $(0,1)$ sur $(3/4,1)$. Ses racines positives d'ordre $30$ récupèrent $q_\pm$. Les identités
\[
 \log s_+=-30(x+y),\qquad \log s_-=-30(x-y)
\]
prouvent \eqref{iii:eq:forma-razon} par différence et somme. Les deux logarithmes
sont négatifs ; leur quotient est positif et leur somme n’est pas nulle. Par conséquent,
\[
 \frac12\log\frac{\log s_+}{\log s_-}
   =\frac12\log\frac{x+y}{x-y}
   =\operatorname{artanh}(y/x)=\eta_{\rm el}.
\]
La relation $Z_{\eta_{\rm el}}/Z_*=e^{-\eta_{\rm el}}$, démontrée dans
\eqref{hmtII:eq:ii-elipse-impedancia}, donne \eqref{iii:eq:forma-LC}.
\end{proof}

La composition conserve l’information nécessaire à chaque reconstruction.
Avec $\hbar$, on retrouve
$a_S=\sqrt{2\hbar}\,e^{\eta_{\rm el}/2}$ et
$b_S=\sqrt{2\hbar}\,e^{-\eta_{\rm el}/2}$ ; avec la coordinatisation
$(\omega,Z_*)$, on retrouve l’inductance et la capacité déjà définies.
L’action, l’horloge et la référence dimensionnelle accompagnent la forme comme
données typées. L’inversion de deux coordonnées adimensionnelles ne
reconstruit ni leurs unités ni toute la mémoire de l’état.

La normalisation intervient avant de résoudre les équations cubiques.
Si l’on reçoit $Z_{\rm vac}$ et $c_{\rm vac}$ dans une coordinatisation dimensionnelle,
les relations \eqref{iii:eq:dimensions} déterminent d’abord
\[
 \widehat Z=\frac{Z_{\rm vac}}{u_Su_Q^{-2}},\qquad
 \widehat c=\frac{c_{\rm vac}}{u_C}.
\]
Ce sont les coordonnées auxquelles s'applique \eqref{iii:eq:forma-inversion}. La référence $Z_*$ appartient à la réalisation LC et est expressément conservée dans \eqref{iii:eq:forma-LC}.

Les deux quotients sont reliés au moyen des canaux, mais évaluent
des fonctions différentes :
\begin{equation}
 \widehat Z=\frac{f(s_-)}{f(s_+)},\qquad
 \frac{Z_{\rm LC}}{Z_*}
       =\sqrt{\frac{\log s_-}{\log s_+}}.
 \label{iii:eq:forma-dos-cocientes}
\end{equation}
La reconstruction précédente fournit l'application entre les deux lectures. En particulier, échanger les feuillets inverse les deux quotients, remplace $\eta_{\rm el}$ par $-\eta_{\rm el}$ et conserve $\widehat c$. Dans la limite symétrique $y=0$, les canaux coïncident et les deux quotients valent $1$. En dehors de ce cas, l'égalité d'origine n'impose pas l'égalité des deux fonctions. La forme d'action est récupérable à partir de la réponse constitutive en conservant l'orientation et la normalisation.
